\section{Limitations}
\label{sec:limitations}
The study has the following limitations, each linked to a direction for future work.
\begin{itemize}
  \item \textbf{One database.} We used OpenAlex, not Scopus or Web of Science. Coverage and metadata quality differ \todo{cite OpenAlex coverage studies after reading them, e.g.\ the reference coverage analysis by Culbert et al.\ (2025), found by search but not yet read}. Future work should repeat the search in Scopus or Web of Science and compare the maps.
  \item \textbf{Incomplete reference data.} Only 278 of 485 articles (57.3\%) have reference lists, which limits coupling and co-citation. Many missing lists come from journals that do not deposit references. A subset or a different source would help.
  \item \textbf{Partial year.} 2026 is incomplete, so yearly counts and trends for 2026 must not be read as a full year.
  \item \textbf{Search sensitivity.} Search is limited to titles and abstracts, and variations of the keywords were not tested. A sensitivity check with alternative terms is needed.
  \item \textbf{Model-assisted screening.} Screening used language models; \todo{report the manual validation and its agreement}. Low-confidence decisions remain.
  \item \textbf{Keywords.} Keywords come from OpenAlex's automatic assignment, not from the authors.
  \item \textbf{Cluster labels and layout.} Cluster membership depends on the algorithm and thresholds; the robustness check shows stable membership but low modularity.
\end{itemize}

\section{Future research avenues}
\todo{Make these specific enough to become study designs. Candidates: (i) a Scopus/Web of Science replication; (ii) a content analysis of the supply chain and inventory articles to test how far results transfer to small businesses; (iii) empirical studies of agent-based LLM systems for inventory tasks that run on local or in-browser models, evaluated on real transaction data.}

\section{Conclusion}
\todo{Close with the contribution, not only a summary. Draft after the discussion is written.}
